% Options for packages loaded elsewhere
\PassOptionsToPackage{unicode}{hyperref}
\PassOptionsToPackage{hyphens}{url}
\PassOptionsToPackage{dvipsnames,svgnames,x11names}{xcolor}
%
\documentclass[
  11pt,
]{article}

\usepackage{amsmath,amssymb}
\usepackage{iftex}
\ifPDFTeX
  \usepackage[T1]{fontenc}
  \usepackage[utf8]{inputenc}
  \usepackage{textcomp} % provide euro and other symbols
\else % if luatex or xetex
  \usepackage{unicode-math}
  \defaultfontfeatures{Scale=MatchLowercase}
  \defaultfontfeatures[\rmfamily]{Ligatures=TeX,Scale=1}
\fi
\usepackage{lmodern}
\ifPDFTeX\else  
    % xetex/luatex font selection
    \setmainfont[]{Arial}
\fi
% Use upquote if available, for straight quotes in verbatim environments
\IfFileExists{upquote.sty}{\usepackage{upquote}}{}
\IfFileExists{microtype.sty}{% use microtype if available
  \usepackage[]{microtype}
  \UseMicrotypeSet[protrusion]{basicmath} % disable protrusion for tt fonts
}{}
\makeatletter
\@ifundefined{KOMAClassName}{% if non-KOMA class
  \IfFileExists{parskip.sty}{%
    \usepackage{parskip}
  }{% else
    \setlength{\parindent}{0pt}
    \setlength{\parskip}{6pt plus 2pt minus 1pt}}
}{% if KOMA class
  \KOMAoptions{parskip=half}}
\makeatother
\usepackage{xcolor}
\usepackage[a4paper, margin=0.8in]{geometry}
\setlength{\emergencystretch}{3em} % prevent overfull lines
\setcounter{secnumdepth}{5}
% Make \paragraph and \subparagraph free-standing
\makeatletter
\ifx\paragraph\undefined\else
  \let\oldparagraph\paragraph
  \renewcommand{\paragraph}{
    \@ifstar
      \xxxParagraphStar
      \xxxParagraphNoStar
  }
  \newcommand{\xxxParagraphStar}[1]{\oldparagraph*{#1}\mbox{}}
  \newcommand{\xxxParagraphNoStar}[1]{\oldparagraph{#1}\mbox{}}
\fi
\ifx\subparagraph\undefined\else
  \let\oldsubparagraph\subparagraph
  \renewcommand{\subparagraph}{
    \@ifstar
      \xxxSubParagraphStar
      \xxxSubParagraphNoStar
  }
  \newcommand{\xxxSubParagraphStar}[1]{\oldsubparagraph*{#1}\mbox{}}
  \newcommand{\xxxSubParagraphNoStar}[1]{\oldsubparagraph{#1}\mbox{}}
\fi
\makeatother

\usepackage{color}
\usepackage{fancyvrb}
\newcommand{\VerbBar}{|}
\newcommand{\VERB}{\Verb[commandchars=\\\{\}]}
\DefineVerbatimEnvironment{Highlighting}{Verbatim}{commandchars=\\\{\}}
% Add ',fontsize=\small' for more characters per line
\usepackage{framed}
\definecolor{shadecolor}{RGB}{248,248,248}
\newenvironment{Shaded}{\begin{snugshade}}{\end{snugshade}}
\newcommand{\AlertTok}[1]{\textcolor[rgb]{0.94,0.16,0.16}{#1}}
\newcommand{\AnnotationTok}[1]{\textcolor[rgb]{0.56,0.35,0.01}{\textbf{\textit{#1}}}}
\newcommand{\AttributeTok}[1]{\textcolor[rgb]{0.13,0.29,0.53}{#1}}
\newcommand{\BaseNTok}[1]{\textcolor[rgb]{0.00,0.00,0.81}{#1}}
\newcommand{\BuiltInTok}[1]{#1}
\newcommand{\CharTok}[1]{\textcolor[rgb]{0.31,0.60,0.02}{#1}}
\newcommand{\CommentTok}[1]{\textcolor[rgb]{0.56,0.35,0.01}{\textit{#1}}}
\newcommand{\CommentVarTok}[1]{\textcolor[rgb]{0.56,0.35,0.01}{\textbf{\textit{#1}}}}
\newcommand{\ConstantTok}[1]{\textcolor[rgb]{0.56,0.35,0.01}{#1}}
\newcommand{\ControlFlowTok}[1]{\textcolor[rgb]{0.13,0.29,0.53}{\textbf{#1}}}
\newcommand{\DataTypeTok}[1]{\textcolor[rgb]{0.13,0.29,0.53}{#1}}
\newcommand{\DecValTok}[1]{\textcolor[rgb]{0.00,0.00,0.81}{#1}}
\newcommand{\DocumentationTok}[1]{\textcolor[rgb]{0.56,0.35,0.01}{\textbf{\textit{#1}}}}
\newcommand{\ErrorTok}[1]{\textcolor[rgb]{0.64,0.00,0.00}{\textbf{#1}}}
\newcommand{\ExtensionTok}[1]{#1}
\newcommand{\FloatTok}[1]{\textcolor[rgb]{0.00,0.00,0.81}{#1}}
\newcommand{\FunctionTok}[1]{\textcolor[rgb]{0.13,0.29,0.53}{\textbf{#1}}}
\newcommand{\ImportTok}[1]{#1}
\newcommand{\InformationTok}[1]{\textcolor[rgb]{0.56,0.35,0.01}{\textbf{\textit{#1}}}}
\newcommand{\KeywordTok}[1]{\textcolor[rgb]{0.13,0.29,0.53}{\textbf{#1}}}
\newcommand{\NormalTok}[1]{#1}
\newcommand{\OperatorTok}[1]{\textcolor[rgb]{0.81,0.36,0.00}{\textbf{#1}}}
\newcommand{\OtherTok}[1]{\textcolor[rgb]{0.56,0.35,0.01}{#1}}
\newcommand{\PreprocessorTok}[1]{\textcolor[rgb]{0.56,0.35,0.01}{\textit{#1}}}
\newcommand{\RegionMarkerTok}[1]{#1}
\newcommand{\SpecialCharTok}[1]{\textcolor[rgb]{0.81,0.36,0.00}{\textbf{#1}}}
\newcommand{\SpecialStringTok}[1]{\textcolor[rgb]{0.31,0.60,0.02}{#1}}
\newcommand{\StringTok}[1]{\textcolor[rgb]{0.31,0.60,0.02}{#1}}
\newcommand{\VariableTok}[1]{\textcolor[rgb]{0.00,0.00,0.00}{#1}}
\newcommand{\VerbatimStringTok}[1]{\textcolor[rgb]{0.31,0.60,0.02}{#1}}
\newcommand{\WarningTok}[1]{\textcolor[rgb]{0.56,0.35,0.01}{\textbf{\textit{#1}}}}

\providecommand{\tightlist}{%
  \setlength{\itemsep}{0pt}\setlength{\parskip}{0pt}}\usepackage{longtable,booktabs,array}
\usepackage{calc} % for calculating minipage widths
% Correct order of tables after \paragraph or \subparagraph
\usepackage{etoolbox}
\makeatletter
\patchcmd\longtable{\par}{\if@noskipsec\mbox{}\fi\par}{}{}
\makeatother
% Allow footnotes in longtable head/foot
\IfFileExists{footnotehyper.sty}{\usepackage{footnotehyper}}{\usepackage{footnote}}
\makesavenoteenv{longtable}
\usepackage{graphicx}
\makeatletter
\newsavebox\pandoc@box
\newcommand*\pandocbounded[1]{% scales image to fit in text height/width
  \sbox\pandoc@box{#1}%
  \Gscale@div\@tempa{\textheight}{\dimexpr\ht\pandoc@box+\dp\pandoc@box\relax}%
  \Gscale@div\@tempb{\linewidth}{\wd\pandoc@box}%
  \ifdim\@tempb\p@<\@tempa\p@\let\@tempa\@tempb\fi% select the smaller of both
  \ifdim\@tempa\p@<\p@\scalebox{\@tempa}{\usebox\pandoc@box}%
  \else\usebox{\pandoc@box}%
  \fi%
}
% Set default figure placement to htbp
\def\fps@figure{htbp}
\makeatother

\makeatletter
\@ifpackageloaded{caption}{}{\usepackage{caption}}
\AtBeginDocument{%
\ifdefined\contentsname
  \renewcommand*\contentsname{Table of contents}
\else
  \newcommand\contentsname{Table of contents}
\fi
\ifdefined\listfigurename
  \renewcommand*\listfigurename{List of Figures}
\else
  \newcommand\listfigurename{List of Figures}
\fi
\ifdefined\listtablename
  \renewcommand*\listtablename{List of Tables}
\else
  \newcommand\listtablename{List of Tables}
\fi
\ifdefined\figurename
  \renewcommand*\figurename{Figure}
\else
  \newcommand\figurename{Figure}
\fi
\ifdefined\tablename
  \renewcommand*\tablename{Table}
\else
  \newcommand\tablename{Table}
\fi
}
\@ifpackageloaded{float}{}{\usepackage{float}}
\floatstyle{ruled}
\@ifundefined{c@chapter}{\newfloat{codelisting}{h}{lop}}{\newfloat{codelisting}{h}{lop}[chapter]}
\floatname{codelisting}{Listing}
\newcommand*\listoflistings{\listof{codelisting}{List of Listings}}
\makeatother
\makeatletter
\makeatother
\makeatletter
\@ifpackageloaded{caption}{}{\usepackage{caption}}
\@ifpackageloaded{subcaption}{}{\usepackage{subcaption}}
\makeatother

\usepackage{bookmark}

\IfFileExists{xurl.sty}{\usepackage{xurl}}{} % add URL line breaks if available
\urlstyle{same} % disable monospaced font for URLs
\hypersetup{
  pdftitle={06\_Transcriptome},
  pdfauthor={Henry Janse van Rensburg},
  colorlinks=true,
  linkcolor={blue},
  filecolor={Maroon},
  citecolor={Blue},
  urlcolor={Blue},
  pdfcreator={LaTeX via pandoc}}


\title{06\_Transcriptome}
\author{Henry Janse van Rensburg}
\date{2025-06-27}

\begin{document}
\maketitle

\renewcommand*\contentsname{Table of contents}
{
\hypersetup{linkcolor=}
\setcounter{tocdepth}{2}
\tableofcontents
}

\section{Background}\label{background}

We conducted this experiment to investigate the effect of the mmf1
mutation on the transcriptome of \emph{Arabidopsis}. The mmf1 mutation
is a deletion of the gene encoding a TNL protein. Plants lacking mmf1
does not show BX-growth feedbacks, and assemble a unique microbiome. We
hypothesize that the deletion of the MMF1 gene will lead to changes in
the expression of genes involved downstream signalling pathways.

\subsection{Materials and methods}\label{materials-and-methods}

We grew wild-type and mmf1 mutant plants under controlled conditions on
soil conditioned in the greenhouse (BS04) for 6 weeks. We harvested the
roots and split them into samples for microbiome analysis and
transcriptome. Three or two roots were pooled per sample. We etracted
RNA and sequenced the RNA using the Illumina platform. We used the
\texttt{DESeq2} package in R to identify differentially expressed genes.

\begin{itemize}
\tightlist
\item
  \textbf{Purpose}: Perform PCA on the transformed data to identify
  major sources of variation among samples. The PCA plot helps visualize
  how different conditions and genotypes cluster together, revealing
  potential patterns in the data. The plot is saved as a high-resolution
  PDF.
\end{itemize}

\subsection{Important results}\label{important-results}

We see a clear seperation between conditions and genotypes in the PCA
plot. The PERMANOVA analysis confirms that both genotype and
conditioning have a significant effect on the transcriptome, especially
on the BXplus microbiome. The interaction between genotype and
conditioning also shows a significant effect, indicating a combined
influence on gene expression. \#\#\# Conclusions

\section{Setup}\label{setup}

\section{Sample exploring}\label{sample-exploring}

\textbf{Method:} We import the (not normalized) gene count table, the
corresponding meta-files and conduct the sequence depth. Then, We run
the standard differential gene expression analysis with DESeq on the
whole data set and perform shrinkage for better visualization. Then we
visualize the samples in a heatmap and a principle component analysis.

\subsection{Reading Data}\label{reading-data}

\begin{itemize}
\tightlist
\item
  \textbf{Purpose}: Load the raw count data (\texttt{cts}) and sample
  metadata (\texttt{meta}) into R. These datasets are essential for
  downstream differential expression analysis.
\end{itemize}

\subsection{Preparing Metadata}\label{preparing-metadata}

\begin{itemize}
\tightlist
\item
  \textbf{Purpose}: Convert the \texttt{genotype} and
  \texttt{conditioning} columns to factors with specific levels,
  ensuring consistent interpretation in the analysis. Also, define
  colors for \texttt{conditioning} levels to use in visualizations.
\end{itemize}

\subsection{Sequencing Depth
Visualization}\label{sequencing-depth-visualization}

\pandocbounded{\includegraphics[keepaspectratio]{06_Transcriptome_files/figure-pdf/Read count-1.pdf}}

\begin{itemize}
\tightlist
\item
  \textbf{Purpose}: Visualize the distribution of sequencing depth
  across genotypes, with jittered points showing individual samples and
  colored by conditioning. This helps check for any biases in sequencing
  depth across groups.
\end{itemize}

\subsection{Creating DESeq2 Object and Running Differential Expression
Analysis}\label{creating-deseq2-object-and-running-differential-expression-analysis}

\begin{itemize}
\tightlist
\item
  \textbf{Purpose}: Set up the DESeq2 object for differential expression
  analysis, filter out low-count genes, relevel factors, and run the
  analysis. The result is a table (\texttt{res}) containing statistics
  like log2 fold changes and p-values for each gene.
\end{itemize}

\subsection{Log Fold Change Shrinkage}\label{log-fold-change-shrinkage}

\begin{longtable}[]{@{}lrrrrr@{}}
\caption{Shrunken log fold changes.}\tabularnewline
\toprule\noalign{}
& baseMean & log2FoldChange & lfcSE & pvalue & padj \\
\midrule\noalign{}
\endfirsthead
\toprule\noalign{}
& baseMean & log2FoldChange & lfcSE & pvalue & padj \\
\midrule\noalign{}
\endhead
\bottomrule\noalign{}
\endlastfoot
AT1G01010 & 1792.0811 & 0.0059300 & 0.0477818 & 0.8593285 & 0.9528780 \\
AT1G01020 & 509.4902 & -0.0240528 & 0.0634101 & 0.4697863 & 0.7667320 \\
AT1G01030 & 126.3681 & -0.0264731 & 0.0696804 & 0.3468553 & 0.6820738 \\
\end{longtable}

\begin{itemize}
\tightlist
\item
  \textbf{Purpose}: Shrink the log fold changes for better
  interpretability, especially for genes with low counts. This reduces
  noise and makes the results more reliable.
\end{itemize}

\subsection{Variance Stabilizing Transformation (VST) and Heatmap
Preparation}\label{variance-stabilizing-transformation-vst-and-heatmap-preparation}

\begin{itemize}
\tightlist
\item
  \textbf{Purpose}: Transform the count data to stabilize variance,
  making it suitable for clustering and heatmaps. Select the top 20 most
  highly expressed genes for focused visualization.
\end{itemize}

\subsection{Sample Clustering Heatmap}\label{sample-clustering-heatmap}

\begin{itemize}
\tightlist
\item
  \textbf{Purpose}: Visualize the similarity between samples using a
  heatmap of Euclidean distances. This helps assess sample clustering,
  which is important for identifying potential batch effects or
  outliers.
\end{itemize}

\subsection{Principal Component Analysis
(PCA)**}\label{principal-component-analysis-pca}

\begin{longtable}[]{@{}lrrllll@{}}
\caption{PCA data.}\tabularnewline
\toprule\noalign{}
& PC1 & PC2 & group & genotype & name & conditioning \\
\midrule\noalign{}
\endfirsthead
\toprule\noalign{}
& PC1 & PC2 & group & genotype & name & conditioning \\
\midrule\noalign{}
\endhead
\bottomrule\noalign{}
\endlastfoot
mmf1m1 & -15.248402 & 2.887772 & mmf1 & mmf1 & mmf1m1 & minus \\
mmf1m2 & -9.221483 & 4.011595 & mmf1 & mmf1 & mmf1m2 & minus \\
mmf1m3 & 0.184031 & 4.420461 & mmf1 & mmf1 & mmf1m3 & minus \\
\end{longtable}

\pandocbounded{\includegraphics[keepaspectratio]{06_Transcriptome_files/figure-pdf/PCA transcriptome-1.pdf}}

\subsection{PERMANOVA Analysis of Lines within Conditioning
Levels}\label{permanova-analysis-of-lines-within-conditioning-levels}

\begin{itemize}
\tightlist
\item
  \textbf{Purpose}: Convert the \texttt{genotype} and
  \texttt{conditioning} columns in \texttt{vsPCA} to factors to ensure
  that the subsequent PERMANOVA analysis correctly interprets these
  variables as categorical.
\end{itemize}

\begin{itemize}
\tightlist
\item
  \textbf{Purpose}: Extract the different levels of
  \texttt{conditioning} and initialize an empty list \texttt{results} to
  store the PERMANOVA results for each conditioning level.
\end{itemize}

\begin{itemize}
\tightlist
\item
  \textbf{Purpose}:

  \begin{itemize}
  \tightlist
  \item
    \textbf{Loop}: For each conditioning level, the code:

    \begin{itemize}
    \tightlist
    \item
      Subsets the data for that conditioning level.
    \item
      Calculates a distance matrix based on the first two principal
      components (\texttt{PC1} and \texttt{PC2}).
    \item
      Performs PERMANOVA (\texttt{adonis}) to test if \texttt{genotype}
      significantly explains differences in the distance matrix.
    \item
      Stores the ANOVA table from the PERMANOVA result in the
      \texttt{results} list.
    \end{itemize}
  \end{itemize}
\end{itemize}

\begin{verbatim}
Conditioning: minus 
Permutation: free
Number of permutations: 999

Terms added sequentially (first to last)

          Df SumsOfSqs MeanSqs F.Model     R2 Pr(>F)   
genotype   1    433.39  433.39  13.119 0.6212   0.01 **
Residuals  8    264.27   33.03         0.3788          
Total      9    697.66                 1.0000          
---
Signif. codes:  0 '***' 0.001 '**' 0.01 '*' 0.05 '.' 0.1 ' ' 1

Conditioning: plus 
Permutation: free
Number of permutations: 999

Terms added sequentially (first to last)

          Df SumsOfSqs MeanSqs F.Model      R2 Pr(>F)  
genotype   1    883.69  883.69  16.054 0.66742  0.015 *
Residuals  8    440.35   55.04         0.33258         
Total      9   1324.04                 1.00000         
---
Signif. codes:  0 '***' 0.001 '**' 0.01 '*' 0.05 '.' 0.1 ' ' 1
\end{verbatim}

\begin{itemize}
\tightlist
\item
  \textbf{Purpose}: Print the PERMANOVA results for each conditioning
  level to examine how \texttt{genotype} influences variation within
  each conditioning level.
\end{itemize}

\subsection{PERMANOVA Analysis of Conditioning within
Genotypes}\label{permanova-analysis-of-conditioning-within-genotypes}

\begin{itemize}
\tightlist
\item
  \textbf{Purpose}: Reaffirm that the \texttt{genotype} and
  \texttt{conditioning} columns are treated as factors in the analysis.
\end{itemize}

\begin{itemize}
\tightlist
\item
  \textbf{Purpose}: Extract the different levels of \texttt{genotype}
  and initialize an empty list \texttt{results} to store the PERMANOVA
  results for each genotype.
\end{itemize}

\begin{itemize}
\tightlist
\item
  \textbf{Purpose}:

  \begin{itemize}
  \tightlist
  \item
    \textbf{Loop}: For each genotype level, the code:

    \begin{itemize}
    \tightlist
    \item
      Subsets the data for that genotype level.
    \item
      Computes a distance matrix based on the first two principal
      components (\texttt{PC1} and \texttt{PC2}).
    \item
      Performs PERMANOVA to test if \texttt{conditioning} significantly
      explains differences in the distance matrix within each genotype.
    \item
      Stores the ANOVA table from the PERMANOVA result in the
      \texttt{results} list.
    \end{itemize}
  \end{itemize}
\end{itemize}

\begin{verbatim}
Genotype: WT 
Permutation: free
Number of permutations: 999

Terms added sequentially (first to last)

             Df SumsOfSqs MeanSqs F.Model      R2 Pr(>F)   
conditioning  1   1131.42 1131.42  31.382 0.79686  0.007 **
Residuals     8    288.43   36.05         0.20314          
Total         9   1419.85                 1.00000          
---
Signif. codes:  0 '***' 0.001 '**' 0.01 '*' 0.05 '.' 0.1 ' ' 1

Genotype: mmf1 
Permutation: free
Number of permutations: 999

Terms added sequentially (first to last)

             Df SumsOfSqs MeanSqs F.Model      R2 Pr(>F)
conditioning  1     50.59  50.593 0.97247 0.10838  0.349
Residuals     8    416.20  52.025         0.89162       
Total         9    466.79                 1.00000       
\end{verbatim}

\begin{itemize}
\tightlist
\item
  \textbf{Purpose}: Print the PERMANOVA results for each genotype to
  evaluate how \texttt{conditioning} influences variation within each
  genotype.
\end{itemize}

\subsection{PERMANOVA Analysis of Lines
vs.~Conditioning}\label{permanova-analysis-of-lines-vs.-conditioning}

\begin{itemize}
\item
  \textbf{Purpose}: Extract the first two principal components
  (\texttt{PC1} and \texttt{PC2}) from \texttt{vsPCA} and compute a
  distance matrix from these scores. This matrix represents the pairwise
  dissimilarities between samples.
\item
  \textbf{Purpose}: Ensure that both \texttt{genotype} and
  \texttt{conditioning} are treated as categorical factors in the
  PERMANOVA analysis.
\end{itemize}

\begin{longtable}[]{@{}
  >{\raggedright\arraybackslash}p{(\linewidth - 12\tabcolsep) * \real{0.3099}}
  >{\raggedleft\arraybackslash}p{(\linewidth - 12\tabcolsep) * \real{0.0423}}
  >{\raggedleft\arraybackslash}p{(\linewidth - 12\tabcolsep) * \real{0.1408}}
  >{\raggedleft\arraybackslash}p{(\linewidth - 12\tabcolsep) * \real{0.1408}}
  >{\raggedleft\arraybackslash}p{(\linewidth - 12\tabcolsep) * \real{0.1268}}
  >{\raggedleft\arraybackslash}p{(\linewidth - 12\tabcolsep) * \real{0.1408}}
  >{\raggedleft\arraybackslash}p{(\linewidth - 12\tabcolsep) * \real{0.0986}}@{}}
\caption{PERMANOVA results.}\tabularnewline
\toprule\noalign{}
\begin{minipage}[b]{\linewidth}\raggedright
\end{minipage} & \begin{minipage}[b]{\linewidth}\raggedleft
Df
\end{minipage} & \begin{minipage}[b]{\linewidth}\raggedleft
SumsOfSqs
\end{minipage} & \begin{minipage}[b]{\linewidth}\raggedleft
MeanSqs
\end{minipage} & \begin{minipage}[b]{\linewidth}\raggedleft
F.Model
\end{minipage} & \begin{minipage}[b]{\linewidth}\raggedleft
R2
\end{minipage} & \begin{minipage}[b]{\linewidth}\raggedleft
Pr(\textgreater F)
\end{minipage} \\
\midrule\noalign{}
\endfirsthead
\toprule\noalign{}
\begin{minipage}[b]{\linewidth}\raggedright
\end{minipage} & \begin{minipage}[b]{\linewidth}\raggedleft
Df
\end{minipage} & \begin{minipage}[b]{\linewidth}\raggedleft
SumsOfSqs
\end{minipage} & \begin{minipage}[b]{\linewidth}\raggedleft
MeanSqs
\end{minipage} & \begin{minipage}[b]{\linewidth}\raggedleft
F.Model
\end{minipage} & \begin{minipage}[b]{\linewidth}\raggedleft
R2
\end{minipage} & \begin{minipage}[b]{\linewidth}\raggedleft
Pr(\textgreater F)
\end{minipage} \\
\midrule\noalign{}
\endhead
\bottomrule\noalign{}
\endlastfoot
genotype & 1 & 962.7357 & 962.73565 & 21.86099 & 0.3378763 & 0.001 \\
conditioning & 1 & 827.6736 & 827.67361 & 18.79412 & 0.2904756 &
0.001 \\
genotype:conditioning & 1 & 354.3407 & 354.34066 & 8.04607 & 0.1243574 &
0.004 \\
Residuals & 16 & 704.6236 & 44.03897 & NA & 0.2472907 & NA \\
Total & 19 & 2849.3735 & NA & NA & 1.0000000 & NA \\
\end{longtable}

\begin{itemize}
\tightlist
\item
  \textbf{Purpose}:

  \begin{itemize}
  \tightlist
  \item
    \textbf{PERMANOVA}: Perform a PERMANOVA to assess whether
    \texttt{genotype}, \texttt{conditioning}, and their interaction
    significantly explain the variation in the distance matrix.
  \item
    \textbf{Interaction}: The \texttt{genotype\ *\ conditioning} formula
    tests both the main effects of \texttt{genotype} and
    \texttt{conditioning} as well as their interaction.
  \item
    \textbf{Output}: Print the ANOVA table, which provides p-values and
    statistics to evaluate the significance of the factors.
  \end{itemize}
\end{itemize}

\section{Differential expression}\label{differential-expression}

\subsection{Setup and Library Loading}\label{setup-and-library-loading}

\subsection{Data Import and
Preprocessing}\label{data-import-and-preprocessing}

\begin{itemize}
\tightlist
\item
  \textbf{Purpose}: Load the RNA-seq count data and associated metadata
  for downstream analysis.Ensure that \texttt{genotype} and
  \texttt{conditioning} are treated as categorical variables with
  predefined levels.Assign specific colors to the conditioning levels
  for consistent plotting.
\end{itemize}

\subsection{Import data}\label{import-data}

\begin{itemize}
\tightlist
\item
  \textbf{Purpose}: Conduct differential expression analysis for mmf1
  mutants and extract significant results, adding Gene IDs for
  reference.
\end{itemize}

\subsection{Plot Creation Function}\label{plot-creation-function}

\begin{itemize}
\tightlist
\item
  \textbf{Purpose}: Create a combined volcano and MA plot for
  differential expression analysis results, displaying differentially
  expressed genes.
\end{itemize}

\subsection{Plotting Differential Expression
Results}\label{plotting-differential-expression-results}

\begin{itemize}
\item
  \textbf{Purpose}: Create and save a volcano plot for WT plants,
  displaying differentially expressed genes.
\item
  \textbf{Purpose}: Create and save a volcano plot for mmf1 mutants,
  displaying differentially expressed genes.
\end{itemize}

\pandocbounded{\includegraphics[keepaspectratio]{06_Transcriptome_files/figure-pdf/Volcano plots-1.pdf}}

\begin{itemize}
\tightlist
\item
  \textbf{Purpose}: Combine the volcano plots for WT and mmf1 into one
  figure and save it as a PDF.
\end{itemize}

\pandocbounded{\includegraphics[keepaspectratio]{06_Transcriptome_files/figure-pdf/VENN diagram upregulated-1.pdf}}

\pandocbounded{\includegraphics[keepaspectratio]{06_Transcriptome_files/figure-pdf/VENN diagram downregulated-1.pdf}}

\subsection{Prepare data to compare WT to
mmf1}\label{prepare-data-to-compare-wt-to-mmf1}

\begin{itemize}
\tightlist
\item
  \textbf{Purpose}: Prepare the data by filtering the metadata and count
  data to compare WT to mmf1 samples conditioned with BXplus.
\end{itemize}

\begin{itemize}
\tightlist
\item
  \textbf{Purpose}: Prepare the data by filtering the metadata and count
  data to compare WT to mmf1 samples conditioned with BXplus and
  BXminus.
\end{itemize}

\section{Differential Expression Analysis between lines within
conditioning
levels}\label{differential-expression-analysis-between-lines-within-conditioning-levels}

\begin{itemize}
\tightlist
\item
  \textbf{Purpose}: Perform differential expression analysis for BXplus
  samples, extract significant results, and add gene IDs for reference.
\end{itemize}

\subsection{Split significant genes by direction
(upregulated/downregulated)}\label{split-significant-genes-by-direction-upregulateddownregulated}

\begin{itemize}
\tightlist
\item
  \textbf{Purpose}: Separate the significant genes from BXplus and
  BXminus samples into upregulated and downregulated categories.
\end{itemize}

\subsection{Create and Plot Volcano and MA plots for BXplus and
BXminus}\label{create-and-plot-volcano-and-ma-plots-for-bxplus-and-bxminus}

\begin{itemize}
\tightlist
\item
  \textbf{Purpose}: Create a function to generate volcano and MA plots
  for the soil analysis, combining the two plots into a single figure.
\end{itemize}

\subsection{Plot Volcano and MA Plots for BXplus
Soil}\label{plot-volcano-and-ma-plots-for-bxplus-soil}

\pandocbounded{\includegraphics[keepaspectratio]{06_Transcriptome_files/figure-pdf/Volcano plus-1.pdf}}

\begin{verbatim}
pdf 
  2 
\end{verbatim}

\begin{itemize}
\tightlist
\item
  \textbf{Purpose}: Generate the volcano and MA plots for BXplus soil
  using the \texttt{create\_plot\_soil} function and save it as a
  high-resolution PDF.
\end{itemize}

\subsubsection{\texorpdfstring{\textbf{Plot Volcano and MA Plots for
BXminus
Soil}}{Plot Volcano and MA Plots for BXminus Soil}}\label{plot-volcano-and-ma-plots-for-bxminus-soil}

\pandocbounded{\includegraphics[keepaspectratio]{06_Transcriptome_files/figure-pdf/Volcanon minus-1.pdf}}

\begin{verbatim}
pdf 
  2 
\end{verbatim}

\begin{itemize}
\tightlist
\item
  \textbf{Purpose}: Generate the volcano and MA plots for BXminus soil
  using the \texttt{create\_plot\_soil} function and save it as a
  high-resolution PDF.
\end{itemize}

\subsection{Combine Volcano Plots for BXplus and BXminus
Soils}\label{combine-volcano-plots-for-bxplus-and-bxminus-soils}

\begin{itemize}
\tightlist
\item
  \textbf{Purpose}: Combine the volcano plots for BXminus and BXplus
  soils into a single plot and save it as a PDF file.
\end{itemize}

\subsection{Extract Gene Lists for Further
Analysis}\label{extract-gene-lists-for-further-analysis}

\begin{itemize}
\tightlist
\item
  \textbf{Purpose}: Extract lists of gene IDs for significantly
  upregulated and downregulated genes in BXplus and BXminus soils for
  further analysis.
\end{itemize}

\subsection{Identify Overlapping Genes between BXplus and BXminus
Soils}\label{identify-overlapping-genes-between-bxplus-and-bxminus-soils}

\begin{itemize}
\tightlist
\item
  \textbf{Purpose}: Identify overlapping genes that are upregulated and
  downregulated in both BXplus and BXminus soils for further analysis.
\end{itemize}

\subsection{Create Lists of Overlapping Genes for
Visualization}\label{create-lists-of-overlapping-genes-for-visualization}

\begin{itemize}
\tightlist
\item
  \textbf{Purpose}: Prepare lists of upregulated and downregulated genes
  in both BXplus and BXminus soils for visualization.
\end{itemize}

\subsection{Visualize Overlapping Genes Using Venn
Diagrams}\label{visualize-overlapping-genes-using-venn-diagrams}

\pandocbounded{\includegraphics[keepaspectratio]{06_Transcriptome_files/figure-pdf/VENN overlap condition-1.pdf}}

\pandocbounded{\includegraphics[keepaspectratio]{06_Transcriptome_files/figure-pdf/VENN overlap condition-2.pdf}}

\begin{itemize}
\tightlist
\item
  \textbf{Purpose}: Create Venn diagrams to visualize the overlap of
  differentially expressed genes between Col-0 and mmf1 genotypes in
  BXplus and BXminus soils.
\end{itemize}

\section{GO term Enrichment Analysis}\label{go-term-enrichment-analysis}

\subsection{Setup and Load Required
Libraries}\label{setup-and-load-required-libraries}

\begin{center}\rule{0.5\linewidth}{0.5pt}\end{center}

\subsection{Co-Expression Analysis
Setup}\label{co-expression-analysis-setup}

\begin{itemize}
\tightlist
\item
  \textbf{Purpose}: The code chunk above sets up the R environment by
  loading the necessary libraries, setting the working directory, and
  importing the count data and metadata for the co-expression analysis.
  It also converts the genotype and conditioning variables to factors
  and defines their levels, as well as assigns colors for plotting
  different conditions.
\end{itemize}

\subsection{Differential Gene Expression
Analysis}\label{differential-gene-expression-analysis}

\begin{itemize}
\tightlist
\item
  \textbf{Purpose}: The code chunk above separates the count data and
  metadata for the WT and mmf1 genotypes and performs differential gene
  expression analysis using DESeq2 for each genotype. It also cleans up
  the environment by removing intermediate data to free up memory.
\end{itemize}

\subsection{Co-Expression Analysis (Coseq)
Setup}\label{co-expression-analysis-coseq-setup}

\begin{itemize}
\tightlist
\item
  \textbf{Result}: We idenntify the optimal number of clusters as 8
  based on the distribution of selected cluster numbers and the Adjusted
  Rand Index (ARI) comparison.
\end{itemize}

\subsubsection{\texorpdfstring{\textbf{Final Coseq Analysis with Optimal
Cluster
Number}}{Final Coseq Analysis with Optimal Cluster Number}}\label{final-coseq-analysis-with-optimal-cluster-number}

Best to not rerun this part, but import previously run file

\begin{itemize}
\tightlist
\item
  \textbf{Purpose}: Perform the final co-expression analysis for the WT
  genotype using the previously determined optimal K value of 8
  clusters.
\end{itemize}

\subsection{Figure 5: Annotated
Heatmap}\label{figure-5-annotated-heatmap}

\begin{longtable}[]{@{}
  >{\raggedright\arraybackslash}p{(\linewidth - 4\tabcolsep) * \real{0.1429}}
  >{\raggedleft\arraybackslash}p{(\linewidth - 4\tabcolsep) * \real{0.1143}}
  >{\raggedright\arraybackslash}p{(\linewidth - 4\tabcolsep) * \real{0.7429}}@{}}
\caption{Annotated GO clusters in WT}\tabularnewline
\toprule\noalign{}
\begin{minipage}[b]{\linewidth}\raggedright
geneID
\end{minipage} & \begin{minipage}[b]{\linewidth}\raggedleft
Cluster
\end{minipage} & \begin{minipage}[b]{\linewidth}\raggedright
Description
\end{minipage} \\
\midrule\noalign{}
\endfirsthead
\toprule\noalign{}
\begin{minipage}[b]{\linewidth}\raggedright
geneID
\end{minipage} & \begin{minipage}[b]{\linewidth}\raggedleft
Cluster
\end{minipage} & \begin{minipage}[b]{\linewidth}\raggedright
Description
\end{minipage} \\
\midrule\noalign{}
\endhead
\bottomrule\noalign{}
\endlastfoot
AT1G01090 & 6 & purine nucleoside bisphosphate biosynthetic process \\
AT1G01090 & 6 & acyl-CoA metabolic process \\
AT1G01090 & 6 & acyl-CoA biosynthetic process \\
\end{longtable}

\begin{itemize}
\tightlist
\item
  \textbf{Purpose}: Generate an annotated heatmap of the co-expression
  clusters for the WT genotype, displaying the expression profiles of
  the clusters along with Gene Ontology (GO) annotations. The heatmap is
  saved as a PDF file for further analysis or visualization.
\end{itemize}

\subsection{Function to create the cluster
heatmap}\label{function-to-create-the-cluster-heatmap}

\begin{Shaded}
\begin{Highlighting}[]
\CommentTok{\# Open PDF device}
\FunctionTok{pdf}\NormalTok{(}\StringTok{"WT\_coseq.annotated.heatmap2\_new.pdf"}\NormalTok{, }\AttributeTok{width =} \DecValTok{14}\NormalTok{, }\AttributeTok{height =} \DecValTok{7}\NormalTok{)}
\CommentTok{\# Purpose: Open a PDF device for plotting, specifying the file path and dimensions (width and height). This allows the annotated heatmap to be saved as a PDF.}

\DocumentationTok{\#\# make annotated heatmap}
\NormalTok{heatmap2 }\OtherTok{\textless{}{-}} \FunctionTok{clusterheatmap3}\NormalTok{(WTCL, df\_WT, WTGO, }\AttributeTok{GOnumber =} \DecValTok{5}\NormalTok{, }\AttributeTok{cluster\_order =}\NormalTok{ T)}
\CommentTok{\# Purpose: Generate an annotated heatmap using a custom \textasciigrave{}clusterheatmap\textasciigrave{} function. The heatmap displays the expression profiles of the clusters along with GO annotations. The top 5 GO terms for each cluster are included, and clusters are ordered (\textasciigrave{}cluster\_order = T\textasciigrave{}).}

\CommentTok{\#Close PDF device}
\FunctionTok{dev.off}\NormalTok{()}
\end{Highlighting}
\end{Shaded}

\begin{verbatim}
pdf 
  2 
\end{verbatim}

\begin{Shaded}
\begin{Highlighting}[]
\CommentTok{\# Purpose: Close the PDF device, finalizing the PDF file that contains the annotated heatmap.}
\FunctionTok{print}\NormalTok{(heatmap2)}
\end{Highlighting}
\end{Shaded}

\pandocbounded{\includegraphics[keepaspectratio]{06_Transcriptome_files/figure-pdf/unnamed-chunk-51-1.pdf}}

\subsection{GO Enrichment Data Preparation and
Filtering}\label{go-enrichment-data-preparation-and-filtering}

\begin{itemize}
\tightlist
\item
  \textbf{Purpose}: Prepare the GO enrichment data for plotting by
  filtering the results, selecting the top 5 GO terms with the lowest
  p.adjust values for each cluster, converting p.adjust to
  -log10(p.adjust), simplifying GO term descriptions, and adding an
  \texttt{Expression} column to categorize clusters as ``Upregulated''
  or ``Downregulated.'' The \texttt{Cluster\_number} column is also
  created for labeling in the plot.
\end{itemize}

Sorting and Ordering Data

\begin{itemize}
\tightlist
\item
  \textbf{Purpose}: Re-sort the data frame by \texttt{Cluster} and
  \texttt{Description} to maintain consistency. Reorder the
  \texttt{Description} factor levels to match the order determined by
  \texttt{GeneRatio} and \texttt{Description}.
\end{itemize}

\subsection{Plotting GO Enrichment
Results}\label{plotting-go-enrichment-results}

\pandocbounded{\includegraphics[keepaspectratio]{06_Transcriptome_files/figure-pdf/Enrichment dotplot-1.pdf}}

\begin{itemize}
\tightlist
\item
  \textbf{Purpose}: Plot the GO enrichment results for the WT genotype,
  displaying clusters along the x-axis and GO terms along the y-axis.
  The size of points represents \texttt{GeneRatio}, and color represents
  \texttt{p.adjust}. The plot is saved as a PDF for further analysis or
  visualization.
\end{itemize}

\subsection{Session Information}\label{session-information}

\begin{verbatim}
R version 4.4.1 (2024-06-14)
Platform: aarch64-apple-darwin20
Running under: macOS 15.5

Matrix products: default
BLAS:   /Library/Frameworks/R.framework/Versions/4.4-arm64/Resources/lib/libRblas.0.dylib 
LAPACK: /Library/Frameworks/R.framework/Versions/4.4-arm64/Resources/lib/libRlapack.dylib;  LAPACK version 3.12.0

locale:
[1] en_US.UTF-8/en_US.UTF-8/en_US.UTF-8/C/en_US.UTF-8/en_US.UTF-8

time zone: Europe/Copenhagen
tzcode source: internal

attached base packages:
 [1] parallel  grid      stats4    stats     graphics  grDevices utils    
 [8] datasets  methods   base     

other attached packages:
 [1] enrichplot_1.26.6           biomaRt_2.58.2             
 [3] varhandle_2.0.6             doBy_4.6.24                
 [5] data.table_1.17.6           simplifyEnrichment_2.0.0   
 [7] reshape2_1.4.4              splitstackshape_1.4.8      
 [9] DOSE_4.0.0                  ggsci_3.2.0                
[11] clusterProfiler_4.14.6      org.At.tair.db_3.20.0      
[13] AnnotationDbi_1.68.0        doParallel_1.0.17          
[15] iterators_1.0.14            foreach_1.5.2              
[17] coseq_1.30.0                BiocManager_1.30.25        
[19] BiocParallel_1.40.0         ggvenn_0.1.10              
[21] cowplot_1.1.3               vegan_2.6-8                
[23] lattice_0.22-6              permute_0.9-7              
[25] openxlsx_4.2.8              rstudioapi_0.17.1          
[27] knitr_1.49                  ggnewscale_0.5.1           
[29] Cairo_1.6-2                 magrittr_2.0.3             
[31] lubridate_1.9.4             forcats_1.0.0              
[33] stringr_1.5.1               purrr_1.0.2                
[35] readr_2.1.5                 tidyr_1.3.1                
[37] tidyverse_2.0.0             ggVennDiagram_1.5.2        
[39] gridExtra_2.3               PoiClaClu_1.0.2.1          
[41] tibble_3.2.1                glmpca_0.2.0               
[43] ComplexHeatmap_2.22.0       ggrepel_0.9.6              
[45] dplyr_1.1.4                 hexbin_1.28.5              
[47] RColorBrewer_1.1-3          vsn_3.74.0                 
[49] viridis_0.6.5               viridisLite_0.4.2          
[51] scales_1.3.0                apeglm_1.28.0              
[53] ggplot2_3.5.1               DESeq2_1.46.0              
[55] SummarizedExperiment_1.36.0 Biobase_2.66.0             
[57] MatrixGenerics_1.18.0       matrixStats_1.5.0          
[59] GenomicRanges_1.58.0        GenomeInfoDb_1.42.1        
[61] IRanges_2.40.1              S4Vectors_0.44.0           
[63] BiocGenerics_0.52.0        

loaded via a namespace (and not attached):
  [1] fs_1.6.5                httr_1.4.7              numDeriv_2016.8-1.1    
  [4] backports_1.5.0         tools_4.4.1             R6_2.5.1               
  [7] lazyeval_0.2.2          mgcv_1.9-1              GetoptLong_1.0.5       
 [10] withr_3.0.2             prettyunits_1.2.0       bayesm_3.1-6           
 [13] preprocessCore_1.68.0   cli_3.6.3               textshaping_0.4.1      
 [16] microbenchmark_1.5.0    labeling_0.4.3          slam_0.1-55            
 [19] mvtnorm_1.3-2           robustbase_0.99-4-1     tm_0.7-15              
 [22] proxy_0.4-27            systemfonts_1.1.0       yulab.utils_0.1.9      
 [25] gson_0.1.0              R.utils_2.12.3          plotrix_3.8-4          
 [28] bbmle_1.0.25.1          limma_3.62.1            RSQLite_2.3.9          
 [31] generics_0.1.3          gridGraphics_0.5-1      shape_1.4.6.1          
 [34] HTSCluster_2.0.11       zip_2.3.1               GO.db_3.20.0           
 [37] Matrix_1.7-0            abind_1.4-8             R.methodsS3_1.8.2      
 [40] lifecycle_1.0.4         scatterplot3d_0.3-44    yaml_2.3.10            
 [43] edgeR_4.4.1             BiocFileCache_2.10.2    qvalue_2.34.0          
 [46] SparseArray_1.6.0       blob_1.2.4              promises_1.3.2         
 [49] crayon_1.5.3            bdsmatrix_1.3-7         ggtangle_0.0.6         
 [52] KEGGREST_1.42.0         HTSFilter_1.46.0        capushe_1.1.2          
 [55] pillar_1.10.1           fgsea_1.32.0            rjson_0.2.23           
 [58] boot_1.3-30             codetools_0.2-20        fastmatch_1.1-6        
 [61] glue_1.8.0              ggfun_0.1.8             vctrs_0.6.5            
 [64] png_0.1-8               treeio_1.30.0           gtable_0.3.6           
 [67] emdbook_1.3.13          cachem_1.1.0            xfun_0.50              
 [70] S4Arrays_1.6.0          mime_0.12               coda_0.19-4.1          
 [73] simona_1.4.0            statmod_1.5.0           nlme_3.1-164           
 [76] ggtree_3.14.0           bit64_4.5.2             filelock_1.0.3         
 [79] progress_1.2.3          tensorA_0.36.2.1        Deriv_4.1.6            
 [82] affyio_1.76.0           colorspace_2.1-1        DBI_1.2.3              
 [85] tidyselect_1.2.1        curl_6.1.0              bit_4.5.0.1            
 [88] compiler_4.4.1          Rmixmod_2.1.10          compositions_2.0-8     
 [91] xml2_1.3.6              NLP_0.3-2               DelayedArray_0.32.0    
 [94] DEoptimR_1.1-3-1        affy_1.84.0             rappdirs_0.3.3         
 [97] digest_0.6.37           rmarkdown_2.29          XVector_0.46.0         
[100] htmltools_0.5.8.1       pkgconfig_2.0.3         dbplyr_2.5.0           
[103] fastmap_1.2.0           rlang_1.1.4             GlobalOptions_0.1.2    
[106] UCSC.utils_1.2.0        shiny_1.10.0            farver_2.1.2           
[109] jsonlite_1.8.9          GOSemSim_2.32.0         R.oo_1.27.0            
[112] GenomeInfoDbData_1.2.13 ggplotify_0.1.2         patchwork_1.3.0        
[115] munsell_0.5.1           Rcpp_1.0.13-1           ape_5.8-1              
[118] stringi_1.8.4           zlibbioc_1.52.0         MASS_7.3-60.2          
[121] plyr_1.8.9              Biostrings_2.74.1       splines_4.4.1          
[124] hms_1.1.3               circlize_0.4.16         locfit_1.5-9.10        
[127] igraph_2.1.2            XML_3.99-0.18           evaluate_1.0.1         
[130] modelr_0.1.11           tzdb_0.4.0              httpuv_1.6.15          
[133] clue_0.3-66             broom_1.0.7             xtable_1.8-4           
[136] e1071_1.7-16            tidytree_0.4.6          later_1.4.1            
[139] class_7.3-22            ragg_1.3.3              Polychrome_1.5.1       
[142] aplot_0.2.4             memoise_2.0.1           cluster_2.1.6          
[145] corrplot_0.95           timechange_0.3.0       
\end{verbatim}

\begin{itemize}
\tightlist
\item
  \textbf{Purpose}: Print the session information, including details
  about the R version, loaded packages, and the system environment. This
  is useful for reproducibility, as it documents the computational
  environment in which the analysis was performed.
\end{itemize}




\end{document}
